\section{Observability-krav}
\label{sec:krav-observability}

Driftspersonell må kunne overvåke og feilsøke VideoAPI uten 
tilgang til kildekoden. Kravene under beskriver logging, 
helsesjekk og metrikker.

\subsection{Logging av drifts- og sikkerhetshendelser (O1.1)}
\label{subsec:krav-o11}

VideoAPI logger sikkerhets- og driftsrelevante hendelser med 
strukturerte log-nivåer (Info, Warning, Error). Følgende 
hendelser logges: autentiseringsforsøk (suksess og feil), 
opprettelse og sletting av videorom med rom-ID og provider, 
start og stopp av videostrømmer, samt feil i kommunikasjon 
med backend-tjenestene. API-nøkler, JWT-tokens og 
pasientinformasjon logges aldri.

Kravet om å ikke logge hemmeligheter er særlig viktig i 
helsedomenet, der sporing av hendelser må balanseres mot 
personvern og pasientsikkerhet.

\subsection{Health check-endepunkter (O1.2)}
\label{subsec:krav-o12}

API-et eksponerer to endepunkter: \texttt{/health/live} 
returnerer HTTP~200 hvis prosessen kjører, og 
\texttt{/health/ready} returnerer HTTP~200 hvis API-et 
kan nå begge backend-tjenestene. Endepunktene krever ikke 
autentisering slik at overvåkingssystemer kan bruke dem 
uten å håndtere tokens.

\subsection{Prometheus-metrikker (O1.3)}
\label{subsec:krav-o13}

API-et eksponerer et \texttt{/metrics}-endepunkt som 
Prometheus skraper. Implementerte metrikker er:

\begin{itemize}
    \item \texttt{hdo\_active\_rooms\_total\{provider\}}: 
          gauge for antall aktive rom per provider
    \item \texttt{http\_requests\_total\{method,endpoint,status\}}: 
          counter for HTTP-forespørsler
    \item \texttt{http\_request\_duration\_seconds}: 
          histogram for responstider
\end{itemize}

Metrikkene visualiseres i et Grafana-dashboard og gir 
driftspersonell løpende oversikt over systemtilstand.